% Document LaTeX généré depuis PDF
\documentclass[12pt]{article}
\usepackage[utf-8]{inputenc}
\usepackage[french]{babel}
\usepackage[margin=1in]{geometry}
\usepackage{fancyheadings}
\usepackage{listings}
\usepackage{xcolor}

% Configuration des listings
\lstset{
  basicstyle=\ttfamily\small,
  breaklines=true,
  breakatwhitespace=true,
  columns=flexible,
  numbers=left,
  numberstyle=\tiny\color{gray},
  keywordstyle=\color{blue},
  commentstyle=\color{gray},
  stringstyle=\color{red},
  stepnumber=1,
  numbersep=5pt,
  backgroundcolor=\color{white},
  showspaces=false,
  showstringspaces=false,
  showtabs=false,
  frame=single,
  rulecolor=\color{black},
  tabsize=2,
  captionpos=b
}

\title{Document converti depuis PDF}
\author{crypto-graph-explorer}
\date{\today}

\begin{document}

\maketitle

\tableofcontents
\newpage

\section{Contenu du document}

% Contenu extrait du PDF
\begin{verbatim}
                       Licence MIAGE
            Rapport de projet Graphes et OpenData
          Analyse du Réseau WETH sur Polygon
             lors du Lundi Noir du 5 août 2024
                    Projet réalisé du XX au XX
                     Membres du groupe
                       NomPrénom 1
                       NomPrénom 2
                       NomPrénom 3
                   Année universitaire 20242025
                Table des matières
                1 Introduction                                                                                  3
                2 Environnement de travail                                                                      4
                3 Description du projet et objectifs                                                            5
                    3.1  Données utilisées . . . . . . . . . . . . . . . . . . . . . . . . . . . . . . . .      5
                    3.2  Questions de recherche . . . . . . . . . . . . . . . . . . . . . . . . . . . . .       5
                    3.3  Prétraitement des données . . . . . . . . . . . . . . . . . . . . . . . . . . .        6
                4 Bibliothèques, Outils et technologies                                                         7
                    4.1  Analyse de graphe      . . . . . . . . . . . . . . . . . . . . . . . . . . . . . . .   7
                    4.2  Visualisation . . . . . . . . . . . . . . . . . . . . . . . . . . . . . . . . . . .    7
                    4.3  Traitement de données . . . . . . . . . . . . . . . . . . . . . . . . . . . . .        7
                    4.4  Rédaction et compilation . . . . . . . . . . . . . . . . . . . . . . . . . . . .       7
                5 Travail réalisé                                                                               8
                6 RQ1  Analyse de la topologie du réseau                                                       9
                    6.1  Méthodologie . . . . . . . . . . . . . . . . . . . . . . . . . . . . . . . . . .       9
                    6.2  Résultats . . . . . . . . . . . . . . . . . . . . . . . . . . . . . . . . . . . . .   10
                    6.3  Interprétation . . . . . . . . . . . . . . . . . . . . . . . . . . . . . . . . . .    11
                7 RQ2  Analyse de la centralité                                                              13
                    7.1  Méthodologie . . . . . . . . . . . . . . . . . . . . . . . . . . . . . . . . . .      13
                    7.2  Résultats . . . . . . . . . . . . . . . . . . . . . . . . . . . . . . . . . . . . .   14
                    7.3  Interprétation . . . . . . . . . . . . . . . . . . . . . . . . . . . . . . . . . .    17
                8 RQ3  Détection de communautés                                                              19
                    8.1  Méthodologie . . . . . . . . . . . . . . . . . . . . . . . . . . . . . . . . . .      19
                    8.2  Résultats . . . . . . . . . . . . . . . . . . . . . . . . . . . . . . . . . . . . .   20
                    8.3  Interprétation . . . . . . . . . . . . . . . . . . . . . . . . . . . . . . . . . .    22
                9 RQ4  Dynamique temporelle                                                                  24
                    9.1  Méthodologie . . . . . . . . . . . . . . . . . . . . . . . . . . . . . . . . . .      24
                    9.2  Résultats . . . . . . . . . . . . . . . . . . . . . . . . . . . . . . . . . . . . .   25
                    9.3  Interprétation . . . . . . . . . . . . . . . . . . . . . . . . . . . . . . . . . .    27
                10 RQ5  Propriétés petit-monde                                                               29
                    10.1 Méthodologie . . . . . . . . . . . . . . . . . . . . . . . . . . . . . . . . . .      29
                    10.2 Résultats . . . . . . . . . . . . . . . . . . . . . . . . . . . . . . . . . . . . .   30
                    10.3 Interprétation . . . . . . . . . . . . . . . . . . . . . . . . . . . . . . . . . .    32
                11 RQ6  Analyse des ux pondérés                                                             34
                    11.1 Méthodologie . . . . . . . . . . . . . . . . . . . . . . . . . . . . . . . . . .      34
                    11.2 Résultats . . . . . . . . . . . . . . . . . . . . . . . . . . . . . . . . . . . . .   35
                    11.3 Interprétation . . . . . . . . . . . . . . . . . . . . . . . . . . . . . . . . . .    38
                12 Dicultés rencontrées                                                                      40
                                                                1
                13 Bilan                                                                                        41
                    13.1 Conclusion . . . . . . . . . . . . . . . . . . . . . . . . . . . . . . . . . . . .      41
                    13.2 Perspectives . . . . . . . . . . . . . . . . . . . . . . . . . . . . . . . . . . .      41
                A Cahier des charges                                                                            45
                    A.1 Contexte . . . . . . . . . . . . . . . . . . . . . . . . . . . . . . . . . . . . .       45
                    A.2 Objectifs . . . . . . . . . . . . . . . . . . . . . . . . . . . . . . . . . . . . .      45
                    A.3 Périmètre fonctionnel . . . . . . . . . . . . . . . . . . . . . . . . . . . . . .        45
                    A.4 Contraintes techniques . . . . . . . . . . . . . . . . . . . . . . . . . . . . .         46
                    A.5 Livrables attendus       . . . . . . . . . . . . . . . . . . . . . . . . . . . . . . .   46
                B Exemple d'exécution du projet                                                                 47
                    B.1 Installation des dépendances . . . . . . . . . . . . . . . . . . . . . . . . . .         47
                    B.2 Exécution des analyses . . . . . . . . . . . . . . . . . . . . . . . . . . . . .         47
                    B.3 Compilation du rapport . . . . . . . . . . . . . . . . . . . . . . . . . . . .           48
                C Manuel utilisateur                                                                            50
                    C.1 Prérequis      . . . . . . . . . . . . . . . . . . . . . . . . . . . . . . . . . . . .   50
                    C.2 Installation     . . . . . . . . . . . . . . . . . . . . . . . . . . . . . . . . . . .   50
                    C.3 Utilisation . . . . . . . . . . . . . . . . . . . . . . . . . . . . . . . . . . . .      50
                    C.4 Structure des chiers . . . . . . . . . . . . . . . . . . . . . . . . . . . . . .        51
                                                                 2
      1  Introduction
        Le 5 août 2024, les marchés nanciers mondiaux se sont eondrés en quelques heures
      unépisode vite surnommé  Lundi Noir  par la presse nancière. Le déclencheur : la
      haussesurprise des taux directeurs annoncée le 31 juillet 2024 par la Banque du Japon [10],
      qui a forcé le dénouement massif des positions de carry trade sur le yen. Le carry trade
      consiste à emprunter dans une devise à faible taux d'intérêt (le yen, maintenu à des taux
      négatifs depuis 2016) pour investir dans des actifs à rendement plus élevé [11]. Quand le
      yen s'est brusquement apprécié, les investisseurs ont dû liquider leurs positions en même
      temps pour couvrir leurs pertes de change, provoquant des ventes forcées sur toutes les
      classes d'actifs. La contagion a atteint les marchés crypto : sur les blockchains de couche 2
      comme Polygon, les protocoles de nance décentralisée (DeFi) ont subi des vagues de
      liquidations et de transferts de panique. Sur la seule journée du 5 août, plus de 401000
      transferts du jeton WETH (Wrapped Ether) ont été enregistrés sur Polygon.
        Ces transferts forment un réseau orienté pondéré : chaque adresse Ethereum est un
      sommet, chaque transfert une arête orientée de l'émetteur vers le récepteur, et le montant
      transféré dénit le poids de l'arête. La théorie des graphes fournit les outils mathématiques
      pour analyser la topologie de ce réseau, identier les acteurs centraux, mettre en évidence
      la structure communautaire et suivre la dynamique temporelle des échanges [1]. Ce type
      d'analyse va au-delà du simple décompte de transactions : il permet de faire apparaître la
      concentration des ux, le rôle des plateformes d'échange décentralisées, et la fragmentation
      du réseau en période de stress.
        Le rapport s'organise autour de six questions de recherche (RQ) :
        1. RQ1  Topologie du réseau : analyser la structure globale à travers la distri-
         bution des degrés, la densité, les composantes connexes et le diamètre du graphe.
        2. RQ2  Centralité : identier les n÷uds les plus inuents par quatre mesures
         de centralité (degré, betweenness, closeness, PageRank) et interpréter leurs rôles.
        3. RQ3  Communautés : détecter la structure communautaire par l'algorithme
         de Louvain, mesurer la modularité et analyser la distribution des tailles de com-
         munautés.
        4. RQ4  Dynamique temporelle : découper la journée en fenêtres horaires et
         suivre l'évolution des métriques du réseau pour identier les pics d'activité.
        5. RQ5  Propriétés petit-monde : évaluer le coecient de clustering et la
         longueur moyenne des chemins, puis comparer au modèle aléatoire d'Erd®s-Rényi.
        6. RQ6  Flux pondérés : analyser la distribution des volumes WETH échangés,
         quantier la concentration par le coecient de Gini et la courbe de Lorenz.
        Après une présentation de l'environnement de travail et du jeu de données, les six
      questions sont traitées successivement (méthodologie, résultats, interprétation). Un bilan
      conclut le rapport et ouvre quelques perspectives.
                         3
               2     Environnement de travail
                   Voici l'environnement matériel et logiciel utilisé pour le projet.
                                     Table 1  Environnement de développement.
                        Outil                     Version       Rôle
                        Système d'exploitation    Windows 11    Plateforme de développement
                        Python                    3.10+         Langage principal d'analyse
                        pip                       24.x          Gestionnaire de paquets
                        Git                       2.x           Contrôle de version
                        VS Code                   1.9x          Éditeur de code
                                                                             A
                        T X Live                  2024          Distribution LT X
                         E                                                     E
                            A
                        pdfLT X                   (inclus)      Compilation du rapport
                              E
                        Biber                     (inclus)      Backend bibliographique
                        latexmk                   (inclus)      Automatisation compilation
                   Python a été choisi pour ses bibliothèques scientiques (NumPy, Pandas, NetworkX,
               Matplotlib). Les dépendances s'installent via pip install -r requirements.txt.
                                                     A                           A            A
                   Le rapport est compilé avec pdfLT X, plus simple que LuaLT X ou XeLT X et su-
                                                       E                           E            E
               sant pour un document en français avec l'encodage T1. T X Live fournit tous les paquets
                                                                           E
               nécessaires, et latexmk automatise les passes de compilation.
                                                             4
              3    Description du projet et objectifs
              3.1   Données utilisées
                 Lejeudedonnéesutilisédanscetteétudeprovientdelatablepubliquecrypto_polygon
              .transactions hébergée sur Google BigQuery [12]. Les transactions ont été ltrées pour
              ne retenir que les transferts du jeton WETH (Wrapped Ether), identié par l'adresse de
              contrat 0x7ceB23fD6bC0adD59E62ac25578270cFf1b9f619 [13], sur la seule journée du 5
              août 2024. Le chier résultant, polygon_weth_crash_2024-08-05.csv, contient 401709
              transactions pour un volume total de 74Mo. Le Tableau 2 résume les caractéristiques
              principales du jeu de données.
                         Table 2  Caractéristiques du jeu de données WETH Polygon.
                       Propriété      Valeur
                       Source         Google BigQuery (crypto_polygon)
                       Token          WETH(Wrapped Ether)
                       Blockchain     Polygon PoS
                       Date           5 août 2024
                       Transactions   401709
                       Colonnes       transaction_hash, chrono, source, target, weight
                       Taille         74Mo
                 ChaqueenregistrementduchierCSVestdécritparcinqcolonnes:transaction_hash,
              l'identiant unique de la transaction sur la blockchain Polygon; chrono, l'horodatage Unix
              (nombre de secondes depuis l'epoch); source, l'adresse Ethereum du portefeuille émet-
              teur; target, l'adresse Ethereum du portefeuille récepteur; et weight, le montant en
              WETHtransféré.
              3.2   Questions de recherche
                 Le projet s'organise autour de six questions de recherche.
              RQ1  Analyse de la topologie du réseau. Examiner la structure topologique
              globale du réseau : distribution des degrés entrants et sortants, densité, composantes
              faiblement et fortement connexes, diamètre et longueur moyenne des chemins. Déterminer
              si le réseau présente des propriétés de type sans échelle, caractéristiques des réseaux où
              quelques n÷uds concentrent un nombre disproportionné de connexions.
              RQ2  Analyse de la centralité. Identier les n÷uds les plus inuents du réseau
              par quatre mesures de centralité : degré, betweenness, closeness et PageRank. Comparer
              les classements obtenus par chaque mesure et interpréter les rôles des acteurs centraux
              dans l'écosystème DeFi (routeurs DEX, ponts inter-chaînes, cores-forts de protocoles).
              RQ3Détection de communautés. Appliquer l'algorithme de Louvain pour révé-
              ler la structure communautaire sous-jacente du réseau. Mesurer la modularité obtenue et
              analyser la distribution des tailles de communautés. Évaluer le degré de compartimenta-
              lisation du réseau et ses implications pour la propagation des chocs nanciers.
                                                       5
           RQ4  Dynamique temporelle. Découper la journée du 5 août 2024 en fenêtres
           horaires et suivre l'évolution des métriques du réseau au cours du temps. Identier les
           pics d'activité en nombre de transactions et en volume WETH, puis les relier au contexte
           nancier mondial (sessions de marché asiatique, européenne et américaine).
           RQ5  Propriétés petit-monde. Évaluer le coecient de clustering et la longueur
           moyenne des chemins du réseau. Comparer ces métriques au modèle aléatoire d'Erd®s-
           Rényi calibré avec les mêmes paramètres (n et p). Quantier le caractère petit-monde par
           le coecient σ de Humphries et Gurney.
           RQ6  Analyse des ux pondérés. Analyser la distribution des volumes WETH
           échangés sur les arêtes du réseau. Quantier la concentration des ux par le coecient de
           Gini et la courbe de Lorenz. Identier les whales (portefeuilles réalisant des transferts de
           très grande valeur) et caractériser leur rôle dans le réseau.
           3.3  Prétraitement des données
              Le prétraitement comprend trois étapes. D'abord, les arêtes multiples entre une même
           paire d'adresses sont agrégées par somme des poids, avec un attribut tx_count conser-
           vant le nombre de transactions. Ensuite, un graphe orienté pondéré (DiGraph) de 12880
           sommets et 39999 arêtes est construit avec NetworkX; une version non orientée (33827
           arêtes) est dérivée pour RQ3 et RQ5. Enn, un ltrage dust à 10−12 WETH supprime
           1598 arêtes (4%) de bruit numérique pour RQ6.
                                             6
            4    Bibliothèques, Outils et technologies
               Les bibliothèques et technologies utilisées sont décrites ci-dessous, organisées par ca-
            tégorie.
            4.1   Analyse de graphe
            NetworkX (v3.4.2). NetworkX [14] fournit les structures DiGraph et Graph utilisées
            dans ce projet, ainsi que les algorithmes de centralité (degré, betweenness, closeness,
            PageRank), de détection de communautés (Louvain, intégré depuis la version 3.2), de
            composantes connexes, de clustering et de plus courts chemins. C'est le socle commun
            aux six questions de recherche.
            4.2   Visualisation
            Matplotlib (v3.9.4). Matplotlib [15] génère les 16 gures PNG à 200 DPI. Deux
            contextes graphiques sont utilisés : fond sombre (#0F172A) pour les graphes de réseau,
            fond clair pour les graphiques analytiques. La palette de couleurs (bleu #2563EB, violet
            #7C3AED, orange #F59E0B) est centralisée dans src/style.py.
            4.3   Traitement de données
            Pandas (v2.2.3). Pandas charge et manipule le CSV de 401709 transactions : conver-
            sion des timestamps Unix en datetime, agrégation des arêtes multiples via groupby, et
            fenêtrage horaire pour l'analyse temporelle (RQ4).
            NumPy (v2.1). NumPy fournit les tableaux et opérations vectorisées pour les cal-
            culs numériques : distributions de degrés, courbes de Lorenz, statistiques descriptives
            (moyennes, médianes, percentiles).
            SciPy (v1.14). SciPy complète NumPy pour les tests statistiques, les comparaisons de
            distributions et les corrélations de rang (Spearman).
            powerlaw (v2.0.0). Labibliothèque powerlaw ajuste des distributions à queue lourde :
            estimation de l'exposant α et du seuil xmin par maximum de vraisemblance, et tests de
            comparaison entre modèles (loi de puissance, lognormale, exponentielle), selon la métho-
            dologie de Clauset et al. [2]. Utilisée en RQ1 et RQ6.
            4.4   Rédaction et compilation
             A              A                                           A
            LT X avec pdfLT X. Le rapport est composé avec le système LT X, compilé par
               E              E                                           E
                         A
            le moteur pdfLT X. Le document utilise la classe article avec 18 paquets organisés
                           E
            selon un ordre de chargement strict : babel pour les conventions typographiques fran-
            çaises, amsmath pour la notation mathématique, booktabs pour les tableaux profession-
            nels, graphicx pour l'inclusion des gures, biblatex pour la gestion bibliographique,
            hyperref pour les liens cliquables et cleveref pour les références croisées intelligentes.
            Labibliographie est gérée par biblatex avec le backend Biber, orant un support Unicode
            complet pour les noms d'auteurs comportant des caractères accentués.
                                                  7
             5    Travail réalisé
                Le projet s'est déroulé en quatre phases : collecte des données, construction du graphe,
             analyses par question de recherche, et rédaction du rapport.
             Collecte et préparation des données. Les 401709 transferts WETH du 5 août 2024
             ont été extraits de la table publique crypto_polygon.transactions sur Google Big-
             Query[12].Lerésultat,stockédanslechierCSVpolygon_weth_crash_2024-08-05.csv
             (74Mo),aétésoumisàunprétraitementd'agrégation:lesarêtesmultiplesentreunemême
             paire d'adresses (source, target) sont fusionnées par somme des poids (volume WETH) et
             comptage des transactions. Un ltrage dust à 10−12 WETH supprime les micro-transferts
             négligeables. Le graphe résultant comprend 12880 sommets et 39999 arêtes orientées
             pondérées.
             Construction du graphe. Le module graph_builder.py sert d'interface entre les
             données brutes et les six analyses. Sa fonction load_and_build() retourne un triplet
             (df, digraph, graph) : le DataFrame Pandas, un graphe orienté pondéré (DiGraph),
             et un graphe non orienté (Graph) pour RQ3 et RQ5. Toutes les analyses travaillent ainsi
             sur le même jeu de données.
             Analyses par question de recherche. Six modules indépendants (src/rq1_*.py à
             src/rq6_*.py) réalisent chacun une analyse et produisent deux gures PNG à 200 DPI
             (16 au total). Le module src/style.py assure la cohérence graphique : fond sombre pour
             les graphes de réseau, fond clair pour les graphiques analytiques. Le script run_all.py
             exécute les six analyses à la suite.
                                                              A                       A
             Rédaction du rapport. Le rapport est rédigé en LT X et compilé avec pdfLT X +
                                                                E                       E
             Biber. Le chier principal main.tex inclut 17 chiers .tex (sections et annexes). La
             bibliographie est scindée en deux : références académiques et webographie, ltrées par
             mot-clé via biblatex.
                                 Table 3  Répartition des tâches du projet
                     Membre      Tâche                         Contribution
                     [Membre 1]  Collecte BigQuery, RQ1, RQ2   Topologie, centralité
                                                    A
                     [Membre 2]  RQ3, RQ4, rapport LT X        Communautés, temporel
                                                      E
                     [Membre 3]  RQ5, RQ6, style graphique     Petit-monde, ux
                                                     8
           6   RQ1  Analyse de la topologie du réseau
              On commence par examiner la structure topologique globale du réseau : distribution
           des degrés, densité, composantes connexes, diamètre et longueur moyenne des chemins.
           Cesmétriques donnent une première image de l'architecture des échanges WETH pendant
           le crash.
           6.1  Méthodologie
              L'analyse topologique repose sur la construction d'un graphe orienté pondéré G =
           (V,E) à partir des 401000 transferts WETH enregistrés le 5 août 2024. Chaque adresse
           Ethereum constitue un sommet, et chaque transfert une arête orientée du portefeuille
           émetteur vers le portefeuille récepteur. Les arêtes multiples entre une même paire de som-
           mets sont agrégées par somme des poids (volume WETH) et comptage des transactions.
              Quatre familles de métriques sont calculées. Premièrement, la densité du réseau,
           dénie pour un graphe orienté comme :
                                      d =    |E|                             (1)
                                         |V | · (|V | − 1)
           où |V| désigne le nombre de sommets et |E| le nombre d'arêtes. Cette métrique quantie
           la proportion de connexions existantes par rapport au nombre maximal de connexions
           possibles. Le degré moyen est donné par :
                                          ¯  |E|
                                          k = |V|                            (2)
              Deuxièmement, les distributions des degrés entrants (kin) et sortants (kout) sont
           analysées séparément. Plutôt qu'un histogramme classique, nous utilisons la fonction de
           répartition complémentaire cumulative (CCDF), dénie par P(X ≥ k), représentée en
           échelle logarithmique sur les deux axes. Cette représentation est préférable pour les dis-
           tributions à queue lourde car elle évite les artefacts liés au choix des intervalles de re-
           groupement et révèle plus clairement la forme de la distribution dans les hautes valeurs
           de degré [2]. L'observation d'une relation linéaire en échelle log-log suggère un compor-
           tement hétérogène caractéristique des réseaux dits  sans échelle , où quelques n÷uds
           concentrent un nombre disproportionné de connexions [3].
              Troisièmement,lescomposantesconnexessontidentiéesselondeuxdénitions.Les
           composantes faiblement connexes (WCC, weakly connected components) ignorent l'orien-
           tation des arêtes : deux sommets appartiennent à la même WCC s'il existe un chemin
           non orienté entre eux. Les composantes fortement connexes (SCC, strongly connected
           components) respectent l'orientation : deux sommets appartiennent à la même SCC si et
           seulement s'il existe un chemin orienté de l'un à l'autre et réciproquement [1]. La taille
           de la composante géante, exprimée en pourcentage du nombre total de sommets, indique
           le degré de connectivité globale du réseau.
              Quatrièmement, le diamètre (plus long des plus courts chemins) et la longueur
           moyenne des chemins sont estimés. Le diamètre est calculé sur la plus grande compo-
           sante fortement connexe, car il n'est déni que sur les graphes connexes. Lorsque cette
           composante dépasse 5000 sommets, l'excentricité est échantillonnée sur 500 n÷uds aléa-
           toires (graine xée à 42 pour la reproductibilité) an de réduire le coût de calcul. La
           longueur moyenne des chemins est estimée sur la plus grande WCC (version non orientée
           du graphe) par échantillonnage de 1000 paires aléatoires de sommets.
                                             9
             6.2   Résultats
             Métriques de base. LeréseauWETHcomprend|V|=12880sommetset|E|=39999
             arêtes orientées agrégées (33827 arêtes non orientées), pour une densité de d = 2,41×10−4.
                              ¯
             Le degré moyen est k = 3,11. Ces valeurs témoignent d'un réseau extrêmement creux : sur
             les quelque 165 millions de connexions possibles, seule une inme fraction est eectivement
             réalisée. La médiane des degrés entrants et sortants est de 1,0, ce qui signie que la
             majorité des adresses n'ont interagi qu'avec un seul autre portefeuille au cours de la
             journée.
             Distribution des degrés.  LaFigure1présentelaCCDFdesdegrésentrantsetsortants
             en échelle logarithmique. Les deux distributions présentent une queue lourde caractéris-
             tique : le degré entrant maximal atteint kmax = 865 tandis que le degré sortant maximal
                                                  in
             culmineàkmax = 2003.LadécroissanceapproximativementlinéairedelaCCDFenéchelle
                       out
             log-log suggère un comportement compatible avec une distribution à queue lourde, typique
             des réseaux sans échelle. Toutefois, nous ne qualions pas formellement cette distribution
             de loi de puissance en l'absence de test statistique dédié (test de Kolmogorov-Smirnov,
             estimation du paramètre α et de xmin) [2]. L'asymétrie entre degrés entrants et sortants
             le maximum sortant étant plus de deux fois supérieur au maximum entrant  in-
             dique la présence de n÷uds diuseurs (routeurs DEX, ponts inter-chaînes) dont l'activité
             d'émission dépasse largement celle de réception.
             Figure 1  Distribution complémentaire cumulative (CCDF) des degrés entrants et sor-
             tants du réseau WETH sur échelle logarithmique. La queue lourde suggère un comporte-
             ment hétérogène typique des réseaux sans échelle [3].
             Composantes connexes. Le réseau contient 95 composantes faiblement connexes et
             7632 composantes fortement connexes. La composante géante faiblement connexe re-
             groupe 12538 sommets, soit 97,3% du réseau : la quasi-totalité des adresses actives sont
             reliées par des chemins non orientés. En revanche, la composante géante fortement connexe
                                                    10
              ne contient que 5206 sommets (40,4%), ce qui révèle une asymétrie directionnelle impor-
              tante : de nombreuses adresses reçoivent ou émettent des WETH sans réciprocité. Le
              Tableau 4 résume ces résultats. La Figure 2 illustre la distribution de la taille des 20
              plus grandes composantes pour chaque type, en échelle logarithmique en raison de la
              dominance de la composante géante.
                         Table 4  Résumé des composantes connexes du réseau WETH.
                    Métrique                         Faiblement connexes   Fortement connexes
                    Nombre de composantes                              95                 7632
                    Composante géante (sommets)                    12538                  5206
                    Composante géante (%)                          97,3%                 40,4%
              Figure2Distributiondelatailledes20plusgrandescomposantesconnexes(faiblement
              et fortement) du réseau WETH, en échelle logarithmique. La composante géante domine
              nettement dans les deux cas.
              Diamètreetlongueurdeschemins. Lediamètreduréseau,estimésurlaplusgrande
              composante fortement connexe (5206 sommets) par échantillonnage de 500 excentricités
              aléatoires, est d'environ D ≈ 14 sauts. La longueur moyenne des chemins, estimée sur
              la plus grande composante faiblement connexe (12538 sommets, graphe non orienté) à
                                                      ¯
              partir de 1000 paires aléatoires, est de ℓ ≈ 3,97 sauts. Malgré les 12880 sommets du
              réseau, deux adresses quelconques de la composante géante sont séparées en moyenne par
              moins de quatre intermédiaires, ce qui témoigne d'une propagation ecace des transferts
              de valeur.
              6.3    Interprétation
                                                    −4
                 La densité très faible (d = 2,41×10   ) est typique des réseaux de transactions block-
              chain, où la plupart des portefeuilles n'interagissent qu'avec une ou deux contreparties. Ce
              réseau creux est pourtant presque entièrement connecté (WCC à 97,3%) : les transferts
              transitent par un petit nombre de n÷uds centraux  routeurs DEX, ponts inter-chaînes,
              contrats de prêt  qui servent de plaques tournantes [1].
                 Laqueuelourdeestcohérenteavecl'attachementpréférentieldeBarabásietAlbert[3]:
              les adresses les plus connectées attirent davantage de nouvelles connexions. Pendant un
                                                        11
      crash, l'eet s'amplie : les protocoles d'échange et de liquidation concentrent un nombre
      croissant de transactions quand les utilisateurs cherchent tous à convertir ou transférer
      leurs WETH en même temps. L'asymétrie entre degré sortant maximal (2003) et entrant
      maximal (865) suggère que certains contrats  routeurs DEX ou bots de liquidation 
      ont diusé des WETH vers un très grand nombre de destinataires.
        L'écart notable entre la composante géante faiblement connexe (97,3%) et la compo-
      sante géante fortement connexe (40,4%) révèle une structure dite  en n÷ud papillon 
      (bow-tie), typique des réseaux orientés de grande taille : un noyau fortement connecté
      de transferts bidirectionnels, entouré de branches entrantes (adresses qui reçoivent sans
      émettre en retour) et de branches sortantes (adresses qui émettent sans recevoir). Cette
      architecture est conforme aux observations faites sur d'autres réseaux de crypto-actifs.
                                         ¯
        Enn, un diamètre court (D ≈ 14) et une longueur moyenne des chemins de ℓ ≈ 3,97
      pour près de 13000 sommets suggèrent déjà un comportement de type  petit monde 
      (testé formellement en RQ5) : n'importe quelle adresse est séparée de n'importe quelle
      autre par moins de quatre intermédiaires, ce qui a pu faciliter la cascade rapide de trans-
      ferts de panique le 5 août.
        La Figure 3 illustre cette architecture en étoile : le sous-graphe des 20 n÷uds les plus
      connectés et de leurs voisins directs montre la centralisation extrême du réseau autour de
      quelques hubs (en orange), chacun irradiant vers de nombreuses adresses périphériques.
      Figure 3  Sous-graphe des 20 hubs principaux (en orange) du réseau WETH et de leurs
      voisins directs. La topologie en étoile conrme la structure hub-and-spoke caractéristique
      des réseaux sans échelle.
                         12
              7    RQ2  Analyse de la centralité
                 Ladistributiondesdegrésàqueuelourde(RQ1)suggèrequecertainsn÷udsconcentrent
              unnombredisproportionnédeconnexions.Onchercheiciàidentiercesn÷udsparquatre
              mesures de centralité [4], et à déterminer si les mêmes adresses dominent systématique-
              mentles quatre classements ou si chaque mesure capture un aspect diérent de l'inuence.
              7.1    Méthodologie
                 Quatre mesures de centralité sont calculées sur le graphe orienté pondéré G = (V,E)
              construit à partir des 12880 sommets et 39999 arêtes agrégées du réseau WETH. Chaque
              mesureéclaireunedimensiondiérentedel'importanced'unn÷ud:lenombredeconnexions
              directes, le rôle d'intermédiaire, la proximité structurelle et l'importance relative dans un
              modèle de marche aléatoire.
              Centralité de degré. La centralité de degré mesure le nombre de connexions directes
              d'un sommet, normalisé par le nombre maximal de connexions possibles. Pour un graphe
              orienté, on distingue la centralité de degré entrant et sortant :
                                               k (v)                  k  (v)
                                      C (v) =    in   ,    C (v)= out                            (3)
                                       in      |V | − 1      out     |V | − 1
              où k (v) et k  (v) désignent respectivement le degré entrant et sortant du sommet v, et
                  in       out
              |V | = 12880 le nombre total de sommets. Un degré entrant élevé identie les adresses re-
              cevant des transferts depuis de nombreuses sources (récepteurs populaires), tandis qu'un
              degré sortant élevé signale les adresses émettant vers de nombreux destinataires (distri-
              buteurs actifs).
              Centralité d'intermédiarité.     La centralité d'intermédiarité (betweenness centrality),
              introduite par Freeman [4], quantie le rôle d'un sommet comme pont entre les autres
              paires de sommets :
                                                         Xσ(v)
                                               C (v) =        st                                 (4)
                                                 B             σ
                                                        s,t∈V   st
                                                        s̸=v̸=t
              où σst est le nombre de plus courts chemins entre s et t, et σst(v) le nombre de ces che-
              mins passant par v. Un n÷ud à forte intermédiarité constitue un goulot d'étranglement
              potentiel : sa suppression allongerait signicativement les chemins entre de nombreuses
              paires de sommets. Le calcul exact de cette mesure est en O(|V| · |E|), ce qui représente
              environ 5 × 108 opérations pour notre graphe. An de maintenir un temps d'exécution
              raisonnable (environ 5 secondes au lieu de plusieurs heures), nous utilisons un échan-
              tillonnage aléatoire de k = 500 sommets pivots avec une graine xée à 42 pour assurer
              la reproductibilité. Cette approximation est standard dans la littérature pour les graphes
              de plus de 5000 sommets [1].
              Centralité de proximité. La centralité de proximité (closeness centrality) mesure la
              proximité moyenne d'un sommet par rapport à tous les autres :
                                                          |V ′| − 1
                                             CC(v) =    X d(v,u)                                 (5)
                                                      u∈V′,u̸=v
                                                        13
             où d(v,u) désigne la longueur du plus court chemin entre v et u, et |V′| le nombre de
             sommets de la composante considérée. Cette mesure n'est dénie que lorsque tous les
             sommets sont atteignables depuis v (i.e., d(v,u) < ∞ pour tout u). Or, le réseau WETH
             présente 95 composantes faiblement connexes (cf. section 6.2), rendant le calcul impossible
             sur le graphe complet. Nous restreignons donc le calcul à la plus grande composante
             faiblement connexe (WCC) en version non orientée, qui contient 12538 sommets, soit
             97,3% du réseau. Un n÷ud à forte proximité est structurellement  au centre  du réseau,
             accessible en peu de sauts depuis n'importe quel autre sommet.
             PageRank. Le PageRank, introduit par Page et al. [5], modélise l'importance d'un
             sommet par une marche aléatoire avec redirection. À chaque pas, un marcheur aléatoire
             suit une arête sortante avec probabilité (1 − α) ou se téléporte vers un sommet aléatoire
             avec probabilité α. Le score PageRank de v correspond à la probabilité stationnaire de
             présence du marcheur sur v. Nous utilisons le facteur d'amortissement standard α = 0,85
             et la variante pondérée : les poids des arêtes (volumes WETH agrégés) inuencent la
             probabilité de transition, de sorte que les n÷uds recevant des transferts de volume élevé
             obtiennent un score supérieur. Contrairement à la simple centralité de degré, le PageRank
             tient compte de la qualité des connexions entrantes : être connecté à un n÷ud lui-même
             inu
\end{verbatim}

\end{document}
